% ============================
%            ABSTRACT
% ============================

% Esta página no muestra número, sin modificar
% el estilo de las páginas posteriores.
\thispagestyle{empty}

% ============================
%      ENCABEZADO INSTITUCIONAL
% ============================

\noindent
\begin{tabularx}{\textwidth}{
    @{}
    >{\centering\arraybackslash}m{3.0cm}
    >{\centering\arraybackslash}X
    >{\centering\arraybackslash}m{3.0cm}
    @{}
    }

    % Logo IPN
    \includegraphics[
        width=2.4cm,
        height=2.8cm,
        keepaspectratio
    ]{Images/logo-ipn.png}

     &

    % Texto institucional
    \begin{minipage}{\linewidth}
        \centering
        \large\bfseries

        INSTITUTO POLITÉCNICO NACIONAL\\[0.2em]

        ESCUELA SUPERIOR DE CÓMPUTO\\[0.2em]

        SUBDIRECCIÓN ACADÉMICA
    \end{minipage}

     &

    % Logo ESCOM
    \includegraphics[
        width=2.8cm,
        height=2.8cm,
        keepaspectratio
    ]{Images/ESCOM.png}
\end{tabularx}

\vspace{0.8em}

% ============================
%       NÚMERO DE TT Y FECHA
% ============================

\noindent
No. de TT: 2027-A178
\hfill
12 de mayo de 2026

\vspace{0.8em}

% ============================
%       TIPO DE DOCUMENTO
% ============================

\begin{center}
    Documento técnico
\end{center}

\vspace{0.5em}

% ============================
%       TÍTULO DEL PROYECTO
% ============================

\begin{center}

    {\Large\bfseries
        Desarrollo de un General Game Player (GGP) utilizando un
        lenguaje de descripción de reglas para juegos de información
        perfecta
        \par}

    \vspace{0.5em}

    2027-A178

\end{center}

\vspace{0.5em}

% ============================
%            AUTORES
% ============================

\begin{center}

    Presentan

    \vspace{0.5em}

    {\large
        Sánchez Catarino Edgar\\[0.5em]
        Linares Villegas Gustavo
    }

\end{center}

\vspace{1.5em}

% ============================
%            DIRECTOR
% ============================

\begin{center}

    Director

    \vspace{0.5em}

    {\large
        Dr. en C. Enrique Escalante Notario
    }

\end{center}

\vspace{1.8em}

% ============================
%             RESUMEN
% ============================

\begin{center}
    \bfseries RESUMEN
\end{center}

\vspace{0.5em}

\noindent
Este Trabajo Terminal propone el diseño e implementación de un General Game
Player (GGP) orientado a juegos discretos, deterministas, por turnos y de
información perfecta. La propuesta busca desacoplar las reglas particulares de
cada juego del motor de decisión del agente mediante un lenguaje de descripción
externo. A partir de dicho lenguaje, el sistema deberá interpretar la estructura
del juego, construir una representación interna, generar movimientos legales,
aplicar transiciones, reconocer estados terminales y seleccionar acciones con
algoritmos generales de búsqueda. El proyecto se valida con Gato como caso
funcional completo y con escenarios controlados de Ajedrez para analizar la
capacidad de generalización del enfoque sin modificar el núcleo algorítmico del
agente.

\vspace{1.5em}

% ============================
%         PALABRAS CLAVE
% ============================

\noindent
\textbf{Palabras clave:}
General Game Playing, lenguaje de descripción de juegos, inteligencia
artificial, juegos de información perfecta, árboles de búsqueda, Monte Carlo
Tree Search.

